\documentclass[11pt,a4paper]{article}
\usepackage[utf8]{inputenc}
\usepackage[T1]{fontenc}
\usepackage[spanish]{babel}
\usepackage{hyperref}
\usepackage{geometry}
\geometry{margin=2.5cm}
\usepackage{parskip}

\title{350 el candado es la combinación del readme y el pegamento}
\author{Alejandro Valenzuela Rivera\\\small Universidad Autónoma de Nuevo León}
\date{\today}
\begin{document}
\maketitle

\begin{quote}
\textbf{Sesgo declarado.} Este trabajo se sostiene desde el sesgo declarado de su autor, y no finge neutralidad. Lo que afirma se puede tumbar, y ese es el punto: el falsificador está escrito, no escondido.
\end{quote}

\section*{Resumen}
El candado no es la prosa de la tesis: es la combinación del readme y el pegamento, porque eso es lo que la academia lee, y el insulto tiene que ser en tiempo real.

\section{Por qué se sostiene}
Lo que la academia lee es el formato; si el formato no está, no lo lee aunque sea cierto.

\section{De dónde viene}
La autoridad: 'eso les mama a la academia y quiero que el insulto sea en tiempo real'.

\section{El caso}
El candado cero: la tesis corta y el detalle duro en el mismo archivo.

\section{Método}
\textquestiondown{}Se puede citar sin preguntarle al autor qué quiso decir?

\section{Límites, y qué lo tumba}
Un candado que necesita a su autor para significar.

\section{Lo que queda abierto}
\textquestiondown{}El candado se lee solo, sin preguntarle al autor qué quiso decir?

\section*{Cómo se leyó esta tesis}
Esto se escribió en el repositorio man-made-empathy, del que este candado es una pieza: candados/04 el lenguaje/350 el candado es la combinación del readme y el pegamento.md. La máquina que lo compuso lo hizo desde este mismo archivo, y no agregó nada que no estuviera escrito. Si algo de esto es falso, se dice en el mismo sitio, y se corrige con fecha.

\end{document}
